\section{Obfuscation Layer}
\label{sec:obfuscation}

\subsection{Control-flow flattening}

Each lifted function is converted into a set of basic blocks with random
non-zero 64-bit state identifiers. The manifest order of the blocks is shuffled
independently of the body order. Every block ends with a flattened transition,
a return, or a trap. The transition carries a state value and a canceling mix:

\begin{align}
  \text{jump}(s, m) &: \quad s' = s + m - m, \label{eq:jump}\\
  \text{jcc}(c, a, b, m) &: \quad s' = c \neq 0 \;?\; a : b,
  \label{eq:jcc}
\end{align}

where $a$ and $b$ are blinded state identifiers and $m$ is a per-build random
mix. The conditional selection is expressed through a runtime register and
canceling arithmetic rather than a native switch table. The goal is not to hide
that a dispatcher exists; it is to remove the direct correspondence between
native block order and bytecode order.

Dispatchers are validated to stay within the block table. An invalid state is a
\code{BadPayload} halt. The flattening does not change the function's observable
behavior: the transition function is a permutation of the original CFG edges.

\subsection{Opaque guards and decoys}

Selected unconditional jumps are replaced by an opaque guard and a decoy block.
The guard uses the invariant

\begin{equation}
  (x \mathbin{\&} 1) \cdot ((x \mathbin{\&} 1) - 1) = 0,
  \label{eq:opaque}
\end{equation}

which is zero for every 64-bit $x$ because one of two consecutive values is even.
The false path executes a decoy block on scratch registers and rejoins the true
successor. The guard never changes live architectural state; a verifier can
establish this by inspecting liveness before accepting the block. Decoys
increase the number of plausible paths a dynamic trace or symbolic engine must
consider without affecting correctness.

\subsection{Mixed Boolean-arithmetic rewriting}

The expression arena applies exact 64-bit wrapping identities:

\begin{align}
  x + y &= (x \oplus y) + 2(x \land y), \label{eq:mba-add}\\
  x - y &= x + \lnot y + 1, \label{eq:mba-sub}\\
  x \oplus y &= (x \lor y) - (x \land y), \label{eq:mba-xor}\\
  x \land y &= (x \lor y) - (x \oplus y), \label{eq:mba-and}\\
  x \lor y &= (x \oplus y) + (x \land y), \label{eq:mba-or}\\
  \lnot x &= x \oplus (2^{64}-1). \label{eq:mba-not}
\end{align}

The rewrite is bounded by \verb|--mba-depth| (0--8, default 3) and
\verb|--mba-budget| (expression nodes, default 512). Every rewrite is checked
for exact equivalence over boundary values and seeded random inputs; an
out-of-budget rewrite is skipped rather than approximated.

Linear identities alone are reducible by modern simplifiers. Sirus therefore
adds one non-linear wrapper per function: a variable-shift bit recombination and
a multiplication by a data-dependent odd value followed by a six-step Newton
modular inverse. The net effect is the identity over 64-bit wrapping arithmetic,
but the expression contains variable shifts and a modular-inversion chain that a
linear simplifier cannot treat as a linear form. This is a cost-asymmetry
mechanism, not a proof of irreducibility. Section~\ref{sec:evaluation} reports
which external tools reduce which forms.

Constants are blinded as two random shares reassembled with XOR:
\begin{equation}
  c = c_0 \oplus c_1, \qquad c_0 \xleftarrow{\$} \{0,1\}^{64}.
  \label{eq:blinding}
\end{equation}
The shares are emitted as separate immediate leaves; no single instruction
contains the plaintext constant.

\subsection{Randomized virtual ISA}

\subsubsection{Semantics and aliases}

The virtual ISA has 48 semantic operations (Appendix~\ref{app:isa}). Each
semantic has four aliases, for a total of 192 assigned opcodes from a shuffled
256-byte pool. Each alias has:

\begin{itemize}
  \item a unique opcode byte;
  \item its own permutation of the canonical operand roles;
  \item a choice of one of two 256-entry register encoding maps;
  \item a flag selecting little- or big-endian immediate bytes.
\end{itemize}

An instruction is encoded as the opcode byte followed by operands in alias order:
a register role becomes one encoded register byte, and an immediate role becomes
eight bytes in the alias byte order. Register bytes are decoded through the
selected map; an invalid encoding is a \code{BadPayload} halt. Alias selection is
random per instruction, so opcode frequency analysis on one block does not
generalize to the next.

\subsubsection{Encrypted decode table}

The runtime needs the semantic, alias, operand order, register map, and
endianness for each opcode. These are packed into a fixed image of

\begin{equation}
  256 + 256 + 2 \cdot 256 + 192 \cdot 8 + 1 = 2561 \text{ bytes}
  \label{eq:table-size}
\end{equation}

and XOR-masked with a per-build ARX stream keyed by a SipHash-derived value. The
image is decrypted into a fixed stack table once per protected invocation. The
table key is present in the generated code, so the table is obfuscated, not
secret from an attacker who executes the VM. The benefit is that the static
binary does not contain a plaintext opcode-to-semantic map.

The build-time encoder and runtime use the same \code{shared/derive.rs} module.
The seed is 256 bits; no alias table is serialized as plaintext outside the
encrypted image. The two register maps are independent permutations with
validity filters; the semantic and alias tables are filtered at use sites.

\subsection{Payload format and encryption}

\subsubsection{Manifest}

The payload begins with an encrypted 8-byte block count followed by encrypted
16-byte manifest entries, one per block:

\begin{center}
\small
\begin{tabular}{@{}lll@{}}
\toprule
Field & Size & Meaning \\
\midrule
state id & 8 bytes & 64-bit block state \\
body offset & 4 bytes & absolute payload offset \\
body length & 4 bytes & encoded body length \\
\bottomrule
\end{tabular}
\end{center}

Entries are shuffled independently of the body order. The count is bounded by
256. The manifest keystream is position-addressed and uses a per-build integer
endianness. A missing state, an out-of-range offset, or a zero-length body is a
\code{BadPayload} halt.

\subsubsection{Rolling body stream}

Each block body is encrypted with a rolling stream seeded by the global key, the
block state, and the absolute body offset:

\begin{align}
  \sigma_0 &= K \oplus \operatorname{rotl}(s,17) \oplus o \oplus D, \label{eq:body-seed}\\
  \mu_i &= \operatorname{mix}\!\left(\sigma_i \oplus \operatorname{rotl}(i,23) \oplus p_{i-1}\right), \label{eq:body-mix}\\
  \sigma_{i+1} &= \operatorname{rotl}(\sigma_i,7) + \operatorname{rotl}(\mu_i,31)
                  \oplus i \oplus p_{i-1}, \label{eq:body-state}\\
  c_i &= p_i \oplus \operatorname{low8}(\mu_i), \label{eq:body-cipher}
\end{align}

where $K$ is the global key, $s$ the block state, $o$ the body offset, $D$ a
domain separator, $p_{-1}=0$, and \code{mix} is a SplitMix64-style ARX avalanche
function. The construction is symmetric: the same equations decrypt. Bytes are
decoded sequentially into locals; no decrypted body buffer is stored.

This stream is an obfuscation stream, not a cryptographic cipher. It is designed
for high-entropy, position-dependent, rolling obfuscation with low per-byte cost.
SipHash-2-4 remains the primitive for manifest encryption, integrity, and key
derivation. The paper does not claim confidentiality against an attacker who can
execute the VM and observe its state.

\subsubsection{Integrity binding}

The runtime computes

\begin{align}
  h &= \operatorname{SipHash24}_{k_h}(\text{encrypted payload}), \label{eq:hash}\\
  K_{\text{eff}} &= a \oplus b \oplus h \oplus h_{\text{expected}}, \label{eq:key}
\end{align}

where $a$ and $b$ are random split shares in the generated source and $k_h$ is
derived from the seed. At build time $h = h_{\text{expected}}$ and
$K_{\text{eff}}$ is the intended key. A modified payload changes $h$ and therefore
$K_{\text{eff}}$; manifest and body decryption then produce invalid structures.
No comparison branch is used, so a naive patch does not reveal an obvious
"integrity check" site.

The binding is 64-bit and the key shares are present in the generated code. It
detects accidental and naive tampering, not an attacker who rewrites the expected
hash, the key shares, and the payload together. It is an obfuscation and
tamper-evidence layer, not a public-key MAC.

\subsection{Masked at-rest state}

Before execution, the VM derives a register mask for every virtual register and a
private-memory keystream from the effective payload key. Reads XOR the mask back
in; writes XOR it out. The initial stack fill is also stored masked. A passive
process dump therefore observes masked values rather than the plaintext
constants used by the protected function. A decoded value can still exist
transiently in a local while one instruction executes; the masking layer is not
an instruction-level tracer defense.

\subsection{Semantic-handler permutation and layout}

The generated engine reorders its top-level semantic match arms with a
seed-derived Fisher--Yates permutation, keeping the wildcard last. The entry
preambles and engine layout are also permuted. LLVM may canonicalize some switch
lowering, so this is layout mutation, not a claim that every build has a distinct
handler implementation.

\subsection{Differential diversity}

Two builds of the same input with different seeds are compared at equal offsets.
On the release tree, the encrypted payload has an equal-byte ratio of 0.0040
(0.40\%) and the full injected VM segment has an equal-byte ratio of 0.0243
(2.43\%). The remaining similarity is the interpreter skeleton, shared runtime
helpers, and canonical instruction sequences. These numbers are diversity
metrics, not semantic security proofs; Section~\ref{sec:security} discusses the
residual.
